
\documentclass[a4paper,12pt]{article}
\usepackage[margin=1cm]{geometry}
\usepackage{tikz}
\usepackage{qrcode}
\usepackage{multicol}
\usepackage{array}
\usepackage{adjustbox}
\pagestyle{empty}
\begin{document}
% --- Bolinha centralizada ---
\newcommand{\marcacao}{%
    \vfill
    \tikz[baseline=(c.base)] \node (c) [circle, draw, minimum width=0.45cm, minimum height=0.45cm] {};
    \vfill
}

% --- Nova coluna centralizada ---
\newcolumntype{C}[1]{>{\centering\arraybackslash}m{#1}}

% --- Cabeçalho ---
\begin{minipage}[t]{0.7\textwidth}
    \textbf{Escola:} CEF04\_PLAN \\
    \textbf{Turma:} 2L - 9º ANO - Matutino \\
    \textbf{Aluno:} ANA CECÍLIA ALVES DE ALMEIDA
\end{minipage}
\begin{minipage}[t]{0.29\textwidth}
    \raggedleft
    \qrcode[height=2.2cm]{ 535448-ANA_CECÍLIA_ALVES_DE_ALMEIDA-2L }
\end{minipage}

\vspace{0.5cm}
\hrule
\vspace{0.5cm}

% --- Gabarito com 25 questões compactado ---
\newcommand{\Gabarito}{
    \renewcommand{\arraystretch}{1.2} % altura das linhas compacta
    \begin{center}
        \begin{tikzpicture}
            \node (table) [inner sep=0pt, anchor=north west] {
                \begin{tabular}{|C{0.6cm}|C{0.8cm}|C{0.8cm}|C{0.8cm}|C{0.8cm}|}
                    \hline
                    \adjustbox{max width=0.5cm}{\textbf{Nº}} & \textbf{A} & \textbf{B} & \textbf{C} & \textbf{D} \\
                    \hline
                    01 & \marcacao & \marcacao & \marcacao & \marcacao \\ \hline
                    02 & \marcacao & \marcacao & \marcacao & \marcacao \\ \hline
                    03 & \marcacao & \marcacao & \marcacao & \marcacao \\ \hline
                    04 & \marcacao & \marcacao & \marcacao & \marcacao \\ \hline
                    05 & \marcacao & \marcacao & \marcacao & \marcacao \\ \hline
                    06 & \marcacao & \marcacao & \marcacao & \marcacao \\ \hline
                    07 & \marcacao & \marcacao & \marcacao & \marcacao \\ \hline
                    08 & \marcacao & \marcacao & \marcacao & \marcacao \\ \hline
                    09 & \marcacao & \marcacao & \marcacao & \marcacao \\ \hline
                    10 & \marcacao & \marcacao & \marcacao & \marcacao \\ \hline
                    11 & \marcacao & \marcacao & \marcacao & \marcacao \\ \hline
                    12 & \marcacao & \marcacao & \marcacao & \marcacao \\ \hline
                    13 & \marcacao & \marcacao & \marcacao & \marcacao \\ \hline
                    14 & \marcacao & \marcacao & \marcacao & \marcacao \\ \hline
                    15 & \marcacao & \marcacao & \marcacao & \marcacao \\ \hline
                    16 & \marcacao & \marcacao & \marcacao & \marcacao \\ \hline
                    17 & \marcacao & \marcacao & \marcacao & \marcacao \\ \hline
                    18 & \marcacao & \marcacao & \marcacao & \marcacao \\ \hline
                    19 & \marcacao & \marcacao & \marcacao & \marcacao \\ \hline
                    20 & \marcacao & \marcacao & \marcacao & \marcacao \\ \hline
                    21 & \marcacao & \marcacao & \marcacao & \marcacao \\ \hline
                    22 & \marcacao & \marcacao & \marcacao & \marcacao \\ \hline
                    23 & \marcacao & \marcacao & \marcacao & \marcacao \\ \hline
                    24 & \marcacao & \marcacao & \marcacao & \marcacao \\ \hline
                    25 & \marcacao & \marcacao & \marcacao & \marcacao \\ \hline
                \end{tabular}
            };
            
            % --- Marcadores OMR ---
            \fill[black] ([shift={(-7mm,7mm)}]table.north west) rectangle ++(5mm,-5mm);
            \fill[black] ([shift={(7mm,7mm)}]table.north east) rectangle ++(-5mm,-5mm);
            \fill[black] ([shift={(-7mm,-7mm)}]table.south west) rectangle ++(5mm,5mm);
            \fill[black] ([shift={(7mm,-7mm)}]table.south east) rectangle ++(-5mm,5mm);
        \end{tikzpicture}
    \end{center}
    \renewcommand{\arraystretch}{1.0}
}

% Chamada do gabarito
\Gabarito

\newpage
% --- Bolinha centralizada ---
\newcommand{\marcacao}{%
    \vfill
    \tikz[baseline=(c.base)] \node (c) [circle, draw, minimum width=0.45cm, minimum height=0.45cm] {};
    \vfill
}

% --- Nova coluna centralizada ---
\newcolumntype{C}[1]{>{\centering\arraybackslash}m{#1}}

% --- Cabeçalho ---
\begin{minipage}[t]{0.7\textwidth}
    \textbf{Escola:} CEF04\_PLAN \\
    \textbf{Turma:} 2L - 9º ANO - Matutino \\
    \textbf{Aluno:} BRUNA KEITY RODRIGUES DE FONTES
\end{minipage}
\begin{minipage}[t]{0.29\textwidth}
    \raggedleft
    \qrcode[height=2.2cm]{ 480291-BRUNA_KEITY_RODRIGUES_DE_FONTES-2L }
\end{minipage}

\vspace{0.5cm}
\hrule
\vspace{0.5cm}

% --- Gabarito com 25 questões compactado ---
\newcommand{\Gabarito}{
    \renewcommand{\arraystretch}{1.2} % altura das linhas compacta
    \begin{center}
        \begin{tikzpicture}
            \node (table) [inner sep=0pt, anchor=north west] {
                \begin{tabular}{|C{0.6cm}|C{0.8cm}|C{0.8cm}|C{0.8cm}|C{0.8cm}|}
                    \hline
                    \adjustbox{max width=0.5cm}{\textbf{Nº}} & \textbf{A} & \textbf{B} & \textbf{C} & \textbf{D} \\
                    \hline
                    01 & \marcacao & \marcacao & \marcacao & \marcacao \\ \hline
                    02 & \marcacao & \marcacao & \marcacao & \marcacao \\ \hline
                    03 & \marcacao & \marcacao & \marcacao & \marcacao \\ \hline
                    04 & \marcacao & \marcacao & \marcacao & \marcacao \\ \hline
                    05 & \marcacao & \marcacao & \marcacao & \marcacao \\ \hline
                    06 & \marcacao & \marcacao & \marcacao & \marcacao \\ \hline
                    07 & \marcacao & \marcacao & \marcacao & \marcacao \\ \hline
                    08 & \marcacao & \marcacao & \marcacao & \marcacao \\ \hline
                    09 & \marcacao & \marcacao & \marcacao & \marcacao \\ \hline
                    10 & \marcacao & \marcacao & \marcacao & \marcacao \\ \hline
                    11 & \marcacao & \marcacao & \marcacao & \marcacao \\ \hline
                    12 & \marcacao & \marcacao & \marcacao & \marcacao \\ \hline
                    13 & \marcacao & \marcacao & \marcacao & \marcacao \\ \hline
                    14 & \marcacao & \marcacao & \marcacao & \marcacao \\ \hline
                    15 & \marcacao & \marcacao & \marcacao & \marcacao \\ \hline
                    16 & \marcacao & \marcacao & \marcacao & \marcacao \\ \hline
                    17 & \marcacao & \marcacao & \marcacao & \marcacao \\ \hline
                    18 & \marcacao & \marcacao & \marcacao & \marcacao \\ \hline
                    19 & \marcacao & \marcacao & \marcacao & \marcacao \\ \hline
                    20 & \marcacao & \marcacao & \marcacao & \marcacao \\ \hline
                    21 & \marcacao & \marcacao & \marcacao & \marcacao \\ \hline
                    22 & \marcacao & \marcacao & \marcacao & \marcacao \\ \hline
                    23 & \marcacao & \marcacao & \marcacao & \marcacao \\ \hline
                    24 & \marcacao & \marcacao & \marcacao & \marcacao \\ \hline
                    25 & \marcacao & \marcacao & \marcacao & \marcacao \\ \hline
                \end{tabular}
            };
            
            % --- Marcadores OMR ---
            \fill[black] ([shift={(-7mm,7mm)}]table.north west) rectangle ++(5mm,-5mm);
            \fill[black] ([shift={(7mm,7mm)}]table.north east) rectangle ++(-5mm,-5mm);
            \fill[black] ([shift={(-7mm,-7mm)}]table.south west) rectangle ++(5mm,5mm);
            \fill[black] ([shift={(7mm,-7mm)}]table.south east) rectangle ++(-5mm,5mm);
        \end{tikzpicture}
    \end{center}
    \renewcommand{\arraystretch}{1.0}
}

% Chamada do gabarito
\Gabarito

\newpage
% --- Bolinha centralizada ---
\newcommand{\marcacao}{%
    \vfill
    \tikz[baseline=(c.base)] \node (c) [circle, draw, minimum width=0.45cm, minimum height=0.45cm] {};
    \vfill
}

% --- Nova coluna centralizada ---
\newcolumntype{C}[1]{>{\centering\arraybackslash}m{#1}}

% --- Cabeçalho ---
\begin{minipage}[t]{0.7\textwidth}
    \textbf{Escola:} CEF04\_PLAN \\
    \textbf{Turma:} 2L - 9º ANO - Matutino \\
    \textbf{Aluno:} BRUNO DE PAULA MAIA
\end{minipage}
\begin{minipage}[t]{0.29\textwidth}
    \raggedleft
    \qrcode[height=2.2cm]{ 591951-BRUNO_DE_PAULA_MAIA-2L }
\end{minipage}

\vspace{0.5cm}
\hrule
\vspace{0.5cm}

% --- Gabarito com 25 questões compactado ---
\newcommand{\Gabarito}{
    \renewcommand{\arraystretch}{1.2} % altura das linhas compacta
    \begin{center}
        \begin{tikzpicture}
            \node (table) [inner sep=0pt, anchor=north west] {
                \begin{tabular}{|C{0.6cm}|C{0.8cm}|C{0.8cm}|C{0.8cm}|C{0.8cm}|}
                    \hline
                    \adjustbox{max width=0.5cm}{\textbf{Nº}} & \textbf{A} & \textbf{B} & \textbf{C} & \textbf{D} \\
                    \hline
                    01 & \marcacao & \marcacao & \marcacao & \marcacao \\ \hline
                    02 & \marcacao & \marcacao & \marcacao & \marcacao \\ \hline
                    03 & \marcacao & \marcacao & \marcacao & \marcacao \\ \hline
                    04 & \marcacao & \marcacao & \marcacao & \marcacao \\ \hline
                    05 & \marcacao & \marcacao & \marcacao & \marcacao \\ \hline
                    06 & \marcacao & \marcacao & \marcacao & \marcacao \\ \hline
                    07 & \marcacao & \marcacao & \marcacao & \marcacao \\ \hline
                    08 & \marcacao & \marcacao & \marcacao & \marcacao \\ \hline
                    09 & \marcacao & \marcacao & \marcacao & \marcacao \\ \hline
                    10 & \marcacao & \marcacao & \marcacao & \marcacao \\ \hline
                    11 & \marcacao & \marcacao & \marcacao & \marcacao \\ \hline
                    12 & \marcacao & \marcacao & \marcacao & \marcacao \\ \hline
                    13 & \marcacao & \marcacao & \marcacao & \marcacao \\ \hline
                    14 & \marcacao & \marcacao & \marcacao & \marcacao \\ \hline
                    15 & \marcacao & \marcacao & \marcacao & \marcacao \\ \hline
                    16 & \marcacao & \marcacao & \marcacao & \marcacao \\ \hline
                    17 & \marcacao & \marcacao & \marcacao & \marcacao \\ \hline
                    18 & \marcacao & \marcacao & \marcacao & \marcacao \\ \hline
                    19 & \marcacao & \marcacao & \marcacao & \marcacao \\ \hline
                    20 & \marcacao & \marcacao & \marcacao & \marcacao \\ \hline
                    21 & \marcacao & \marcacao & \marcacao & \marcacao \\ \hline
                    22 & \marcacao & \marcacao & \marcacao & \marcacao \\ \hline
                    23 & \marcacao & \marcacao & \marcacao & \marcacao \\ \hline
                    24 & \marcacao & \marcacao & \marcacao & \marcacao \\ \hline
                    25 & \marcacao & \marcacao & \marcacao & \marcacao \\ \hline
                \end{tabular}
            };
            
            % --- Marcadores OMR ---
            \fill[black] ([shift={(-7mm,7mm)}]table.north west) rectangle ++(5mm,-5mm);
            \fill[black] ([shift={(7mm,7mm)}]table.north east) rectangle ++(-5mm,-5mm);
            \fill[black] ([shift={(-7mm,-7mm)}]table.south west) rectangle ++(5mm,5mm);
            \fill[black] ([shift={(7mm,-7mm)}]table.south east) rectangle ++(-5mm,5mm);
        \end{tikzpicture}
    \end{center}
    \renewcommand{\arraystretch}{1.0}
}

% Chamada do gabarito
\Gabarito

\newpage
% --- Bolinha centralizada ---
\newcommand{\marcacao}{%
    \vfill
    \tikz[baseline=(c.base)] \node (c) [circle, draw, minimum width=0.45cm, minimum height=0.45cm] {};
    \vfill
}

% --- Nova coluna centralizada ---
\newcolumntype{C}[1]{>{\centering\arraybackslash}m{#1}}

% --- Cabeçalho ---
\begin{minipage}[t]{0.7\textwidth}
    \textbf{Escola:} CEF04\_PLAN \\
    \textbf{Turma:} 2L - 9º ANO - Matutino \\
    \textbf{Aluno:} DAVID MARTINS E SILVA
\end{minipage}
\begin{minipage}[t]{0.29\textwidth}
    \raggedleft
    \qrcode[height=2.2cm]{ 493740-DAVID_MARTINS_E_SILVA-2L }
\end{minipage}

\vspace{0.5cm}
\hrule
\vspace{0.5cm}

% --- Gabarito com 25 questões compactado ---
\newcommand{\Gabarito}{
    \renewcommand{\arraystretch}{1.2} % altura das linhas compacta
    \begin{center}
        \begin{tikzpicture}
            \node (table) [inner sep=0pt, anchor=north west] {
                \begin{tabular}{|C{0.6cm}|C{0.8cm}|C{0.8cm}|C{0.8cm}|C{0.8cm}|}
                    \hline
                    \adjustbox{max width=0.5cm}{\textbf{Nº}} & \textbf{A} & \textbf{B} & \textbf{C} & \textbf{D} \\
                    \hline
                    01 & \marcacao & \marcacao & \marcacao & \marcacao \\ \hline
                    02 & \marcacao & \marcacao & \marcacao & \marcacao \\ \hline
                    03 & \marcacao & \marcacao & \marcacao & \marcacao \\ \hline
                    04 & \marcacao & \marcacao & \marcacao & \marcacao \\ \hline
                    05 & \marcacao & \marcacao & \marcacao & \marcacao \\ \hline
                    06 & \marcacao & \marcacao & \marcacao & \marcacao \\ \hline
                    07 & \marcacao & \marcacao & \marcacao & \marcacao \\ \hline
                    08 & \marcacao & \marcacao & \marcacao & \marcacao \\ \hline
                    09 & \marcacao & \marcacao & \marcacao & \marcacao \\ \hline
                    10 & \marcacao & \marcacao & \marcacao & \marcacao \\ \hline
                    11 & \marcacao & \marcacao & \marcacao & \marcacao \\ \hline
                    12 & \marcacao & \marcacao & \marcacao & \marcacao \\ \hline
                    13 & \marcacao & \marcacao & \marcacao & \marcacao \\ \hline
                    14 & \marcacao & \marcacao & \marcacao & \marcacao \\ \hline
                    15 & \marcacao & \marcacao & \marcacao & \marcacao \\ \hline
                    16 & \marcacao & \marcacao & \marcacao & \marcacao \\ \hline
                    17 & \marcacao & \marcacao & \marcacao & \marcacao \\ \hline
                    18 & \marcacao & \marcacao & \marcacao & \marcacao \\ \hline
                    19 & \marcacao & \marcacao & \marcacao & \marcacao \\ \hline
                    20 & \marcacao & \marcacao & \marcacao & \marcacao \\ \hline
                    21 & \marcacao & \marcacao & \marcacao & \marcacao \\ \hline
                    22 & \marcacao & \marcacao & \marcacao & \marcacao \\ \hline
                    23 & \marcacao & \marcacao & \marcacao & \marcacao \\ \hline
                    24 & \marcacao & \marcacao & \marcacao & \marcacao \\ \hline
                    25 & \marcacao & \marcacao & \marcacao & \marcacao \\ \hline
                \end{tabular}
            };
            
            % --- Marcadores OMR ---
            \fill[black] ([shift={(-7mm,7mm)}]table.north west) rectangle ++(5mm,-5mm);
            \fill[black] ([shift={(7mm,7mm)}]table.north east) rectangle ++(-5mm,-5mm);
            \fill[black] ([shift={(-7mm,-7mm)}]table.south west) rectangle ++(5mm,5mm);
            \fill[black] ([shift={(7mm,-7mm)}]table.south east) rectangle ++(-5mm,5mm);
        \end{tikzpicture}
    \end{center}
    \renewcommand{\arraystretch}{1.0}
}

% Chamada do gabarito
\Gabarito

\newpage
% --- Bolinha centralizada ---
\newcommand{\marcacao}{%
    \vfill
    \tikz[baseline=(c.base)] \node (c) [circle, draw, minimum width=0.45cm, minimum height=0.45cm] {};
    \vfill
}

% --- Nova coluna centralizada ---
\newcolumntype{C}[1]{>{\centering\arraybackslash}m{#1}}

% --- Cabeçalho ---
\begin{minipage}[t]{0.7\textwidth}
    \textbf{Escola:} CEF04\_PLAN \\
    \textbf{Turma:} 2L - 9º ANO - Matutino \\
    \textbf{Aluno:} ELOISA AZEVEDO MEIRELES
\end{minipage}
\begin{minipage}[t]{0.29\textwidth}
    \raggedleft
    \qrcode[height=2.2cm]{ 531472-ELOISA_AZEVEDO_MEIRELES-2L }
\end{minipage}

\vspace{0.5cm}
\hrule
\vspace{0.5cm}

% --- Gabarito com 25 questões compactado ---
\newcommand{\Gabarito}{
    \renewcommand{\arraystretch}{1.2} % altura das linhas compacta
    \begin{center}
        \begin{tikzpicture}
            \node (table) [inner sep=0pt, anchor=north west] {
                \begin{tabular}{|C{0.6cm}|C{0.8cm}|C{0.8cm}|C{0.8cm}|C{0.8cm}|}
                    \hline
                    \adjustbox{max width=0.5cm}{\textbf{Nº}} & \textbf{A} & \textbf{B} & \textbf{C} & \textbf{D} \\
                    \hline
                    01 & \marcacao & \marcacao & \marcacao & \marcacao \\ \hline
                    02 & \marcacao & \marcacao & \marcacao & \marcacao \\ \hline
                    03 & \marcacao & \marcacao & \marcacao & \marcacao \\ \hline
                    04 & \marcacao & \marcacao & \marcacao & \marcacao \\ \hline
                    05 & \marcacao & \marcacao & \marcacao & \marcacao \\ \hline
                    06 & \marcacao & \marcacao & \marcacao & \marcacao \\ \hline
                    07 & \marcacao & \marcacao & \marcacao & \marcacao \\ \hline
                    08 & \marcacao & \marcacao & \marcacao & \marcacao \\ \hline
                    09 & \marcacao & \marcacao & \marcacao & \marcacao \\ \hline
                    10 & \marcacao & \marcacao & \marcacao & \marcacao \\ \hline
                    11 & \marcacao & \marcacao & \marcacao & \marcacao \\ \hline
                    12 & \marcacao & \marcacao & \marcacao & \marcacao \\ \hline
                    13 & \marcacao & \marcacao & \marcacao & \marcacao \\ \hline
                    14 & \marcacao & \marcacao & \marcacao & \marcacao \\ \hline
                    15 & \marcacao & \marcacao & \marcacao & \marcacao \\ \hline
                    16 & \marcacao & \marcacao & \marcacao & \marcacao \\ \hline
                    17 & \marcacao & \marcacao & \marcacao & \marcacao \\ \hline
                    18 & \marcacao & \marcacao & \marcacao & \marcacao \\ \hline
                    19 & \marcacao & \marcacao & \marcacao & \marcacao \\ \hline
                    20 & \marcacao & \marcacao & \marcacao & \marcacao \\ \hline
                    21 & \marcacao & \marcacao & \marcacao & \marcacao \\ \hline
                    22 & \marcacao & \marcacao & \marcacao & \marcacao \\ \hline
                    23 & \marcacao & \marcacao & \marcacao & \marcacao \\ \hline
                    24 & \marcacao & \marcacao & \marcacao & \marcacao \\ \hline
                    25 & \marcacao & \marcacao & \marcacao & \marcacao \\ \hline
                \end{tabular}
            };
            
            % --- Marcadores OMR ---
            \fill[black] ([shift={(-7mm,7mm)}]table.north west) rectangle ++(5mm,-5mm);
            \fill[black] ([shift={(7mm,7mm)}]table.north east) rectangle ++(-5mm,-5mm);
            \fill[black] ([shift={(-7mm,-7mm)}]table.south west) rectangle ++(5mm,5mm);
            \fill[black] ([shift={(7mm,-7mm)}]table.south east) rectangle ++(-5mm,5mm);
        \end{tikzpicture}
    \end{center}
    \renewcommand{\arraystretch}{1.0}
}

% Chamada do gabarito
\Gabarito

\newpage
% --- Bolinha centralizada ---
\newcommand{\marcacao}{%
    \vfill
    \tikz[baseline=(c.base)] \node (c) [circle, draw, minimum width=0.45cm, minimum height=0.45cm] {};
    \vfill
}

% --- Nova coluna centralizada ---
\newcolumntype{C}[1]{>{\centering\arraybackslash}m{#1}}

% --- Cabeçalho ---
\begin{minipage}[t]{0.7\textwidth}
    \textbf{Escola:} CEF04\_PLAN \\
    \textbf{Turma:} 2L - 9º ANO - Matutino \\
    \textbf{Aluno:} ERIK DE OLIVEIRA VALADARES
\end{minipage}
\begin{minipage}[t]{0.29\textwidth}
    \raggedleft
    \qrcode[height=2.2cm]{ 472317-ERIK_DE_OLIVEIRA_VALADARES-2L }
\end{minipage}

\vspace{0.5cm}
\hrule
\vspace{0.5cm}

% --- Gabarito com 25 questões compactado ---
\newcommand{\Gabarito}{
    \renewcommand{\arraystretch}{1.2} % altura das linhas compacta
    \begin{center}
        \begin{tikzpicture}
            \node (table) [inner sep=0pt, anchor=north west] {
                \begin{tabular}{|C{0.6cm}|C{0.8cm}|C{0.8cm}|C{0.8cm}|C{0.8cm}|}
                    \hline
                    \adjustbox{max width=0.5cm}{\textbf{Nº}} & \textbf{A} & \textbf{B} & \textbf{C} & \textbf{D} \\
                    \hline
                    01 & \marcacao & \marcacao & \marcacao & \marcacao \\ \hline
                    02 & \marcacao & \marcacao & \marcacao & \marcacao \\ \hline
                    03 & \marcacao & \marcacao & \marcacao & \marcacao \\ \hline
                    04 & \marcacao & \marcacao & \marcacao & \marcacao \\ \hline
                    05 & \marcacao & \marcacao & \marcacao & \marcacao \\ \hline
                    06 & \marcacao & \marcacao & \marcacao & \marcacao \\ \hline
                    07 & \marcacao & \marcacao & \marcacao & \marcacao \\ \hline
                    08 & \marcacao & \marcacao & \marcacao & \marcacao \\ \hline
                    09 & \marcacao & \marcacao & \marcacao & \marcacao \\ \hline
                    10 & \marcacao & \marcacao & \marcacao & \marcacao \\ \hline
                    11 & \marcacao & \marcacao & \marcacao & \marcacao \\ \hline
                    12 & \marcacao & \marcacao & \marcacao & \marcacao \\ \hline
                    13 & \marcacao & \marcacao & \marcacao & \marcacao \\ \hline
                    14 & \marcacao & \marcacao & \marcacao & \marcacao \\ \hline
                    15 & \marcacao & \marcacao & \marcacao & \marcacao \\ \hline
                    16 & \marcacao & \marcacao & \marcacao & \marcacao \\ \hline
                    17 & \marcacao & \marcacao & \marcacao & \marcacao \\ \hline
                    18 & \marcacao & \marcacao & \marcacao & \marcacao \\ \hline
                    19 & \marcacao & \marcacao & \marcacao & \marcacao \\ \hline
                    20 & \marcacao & \marcacao & \marcacao & \marcacao \\ \hline
                    21 & \marcacao & \marcacao & \marcacao & \marcacao \\ \hline
                    22 & \marcacao & \marcacao & \marcacao & \marcacao \\ \hline
                    23 & \marcacao & \marcacao & \marcacao & \marcacao \\ \hline
                    24 & \marcacao & \marcacao & \marcacao & \marcacao \\ \hline
                    25 & \marcacao & \marcacao & \marcacao & \marcacao \\ \hline
                \end{tabular}
            };
            
            % --- Marcadores OMR ---
            \fill[black] ([shift={(-7mm,7mm)}]table.north west) rectangle ++(5mm,-5mm);
            \fill[black] ([shift={(7mm,7mm)}]table.north east) rectangle ++(-5mm,-5mm);
            \fill[black] ([shift={(-7mm,-7mm)}]table.south west) rectangle ++(5mm,5mm);
            \fill[black] ([shift={(7mm,-7mm)}]table.south east) rectangle ++(-5mm,5mm);
        \end{tikzpicture}
    \end{center}
    \renewcommand{\arraystretch}{1.0}
}

% Chamada do gabarito
\Gabarito

\newpage
% --- Bolinha centralizada ---
\newcommand{\marcacao}{%
    \vfill
    \tikz[baseline=(c.base)] \node (c) [circle, draw, minimum width=0.45cm, minimum height=0.45cm] {};
    \vfill
}

% --- Nova coluna centralizada ---
\newcolumntype{C}[1]{>{\centering\arraybackslash}m{#1}}

% --- Cabeçalho ---
\begin{minipage}[t]{0.7\textwidth}
    \textbf{Escola:} CEF04\_PLAN \\
    \textbf{Turma:} 2L - 9º ANO - Matutino \\
    \textbf{Aluno:} FILIPE GABRIEL PEREIRA PARDINHO
\end{minipage}
\begin{minipage}[t]{0.29\textwidth}
    \raggedleft
    \qrcode[height=2.2cm]{ 581974-FILIPE_GABRIEL_PEREIRA_PARDINHO-2L }
\end{minipage}

\vspace{0.5cm}
\hrule
\vspace{0.5cm}

% --- Gabarito com 25 questões compactado ---
\newcommand{\Gabarito}{
    \renewcommand{\arraystretch}{1.2} % altura das linhas compacta
    \begin{center}
        \begin{tikzpicture}
            \node (table) [inner sep=0pt, anchor=north west] {
                \begin{tabular}{|C{0.6cm}|C{0.8cm}|C{0.8cm}|C{0.8cm}|C{0.8cm}|}
                    \hline
                    \adjustbox{max width=0.5cm}{\textbf{Nº}} & \textbf{A} & \textbf{B} & \textbf{C} & \textbf{D} \\
                    \hline
                    01 & \marcacao & \marcacao & \marcacao & \marcacao \\ \hline
                    02 & \marcacao & \marcacao & \marcacao & \marcacao \\ \hline
                    03 & \marcacao & \marcacao & \marcacao & \marcacao \\ \hline
                    04 & \marcacao & \marcacao & \marcacao & \marcacao \\ \hline
                    05 & \marcacao & \marcacao & \marcacao & \marcacao \\ \hline
                    06 & \marcacao & \marcacao & \marcacao & \marcacao \\ \hline
                    07 & \marcacao & \marcacao & \marcacao & \marcacao \\ \hline
                    08 & \marcacao & \marcacao & \marcacao & \marcacao \\ \hline
                    09 & \marcacao & \marcacao & \marcacao & \marcacao \\ \hline
                    10 & \marcacao & \marcacao & \marcacao & \marcacao \\ \hline
                    11 & \marcacao & \marcacao & \marcacao & \marcacao \\ \hline
                    12 & \marcacao & \marcacao & \marcacao & \marcacao \\ \hline
                    13 & \marcacao & \marcacao & \marcacao & \marcacao \\ \hline
                    14 & \marcacao & \marcacao & \marcacao & \marcacao \\ \hline
                    15 & \marcacao & \marcacao & \marcacao & \marcacao \\ \hline
                    16 & \marcacao & \marcacao & \marcacao & \marcacao \\ \hline
                    17 & \marcacao & \marcacao & \marcacao & \marcacao \\ \hline
                    18 & \marcacao & \marcacao & \marcacao & \marcacao \\ \hline
                    19 & \marcacao & \marcacao & \marcacao & \marcacao \\ \hline
                    20 & \marcacao & \marcacao & \marcacao & \marcacao \\ \hline
                    21 & \marcacao & \marcacao & \marcacao & \marcacao \\ \hline
                    22 & \marcacao & \marcacao & \marcacao & \marcacao \\ \hline
                    23 & \marcacao & \marcacao & \marcacao & \marcacao \\ \hline
                    24 & \marcacao & \marcacao & \marcacao & \marcacao \\ \hline
                    25 & \marcacao & \marcacao & \marcacao & \marcacao \\ \hline
                \end{tabular}
            };
            
            % --- Marcadores OMR ---
            \fill[black] ([shift={(-7mm,7mm)}]table.north west) rectangle ++(5mm,-5mm);
            \fill[black] ([shift={(7mm,7mm)}]table.north east) rectangle ++(-5mm,-5mm);
            \fill[black] ([shift={(-7mm,-7mm)}]table.south west) rectangle ++(5mm,5mm);
            \fill[black] ([shift={(7mm,-7mm)}]table.south east) rectangle ++(-5mm,5mm);
        \end{tikzpicture}
    \end{center}
    \renewcommand{\arraystretch}{1.0}
}

% Chamada do gabarito
\Gabarito

\newpage
% --- Bolinha centralizada ---
\newcommand{\marcacao}{%
    \vfill
    \tikz[baseline=(c.base)] \node (c) [circle, draw, minimum width=0.45cm, minimum height=0.45cm] {};
    \vfill
}

% --- Nova coluna centralizada ---
\newcolumntype{C}[1]{>{\centering\arraybackslash}m{#1}}

% --- Cabeçalho ---
\begin{minipage}[t]{0.7\textwidth}
    \textbf{Escola:} CEF04\_PLAN \\
    \textbf{Turma:} 2L - 9º ANO - Matutino \\
    \textbf{Aluno:} GABRIELLY BEATRYZ MAURÍCIO GOMES
\end{minipage}
\begin{minipage}[t]{0.29\textwidth}
    \raggedleft
    \qrcode[height=2.2cm]{ 570690-GABRIELLY_BEATRYZ_MAURÍCIO_GOMES-2L }
\end{minipage}

\vspace{0.5cm}
\hrule
\vspace{0.5cm}

% --- Gabarito com 25 questões compactado ---
\newcommand{\Gabarito}{
    \renewcommand{\arraystretch}{1.2} % altura das linhas compacta
    \begin{center}
        \begin{tikzpicture}
            \node (table) [inner sep=0pt, anchor=north west] {
                \begin{tabular}{|C{0.6cm}|C{0.8cm}|C{0.8cm}|C{0.8cm}|C{0.8cm}|}
                    \hline
                    \adjustbox{max width=0.5cm}{\textbf{Nº}} & \textbf{A} & \textbf{B} & \textbf{C} & \textbf{D} \\
                    \hline
                    01 & \marcacao & \marcacao & \marcacao & \marcacao \\ \hline
                    02 & \marcacao & \marcacao & \marcacao & \marcacao \\ \hline
                    03 & \marcacao & \marcacao & \marcacao & \marcacao \\ \hline
                    04 & \marcacao & \marcacao & \marcacao & \marcacao \\ \hline
                    05 & \marcacao & \marcacao & \marcacao & \marcacao \\ \hline
                    06 & \marcacao & \marcacao & \marcacao & \marcacao \\ \hline
                    07 & \marcacao & \marcacao & \marcacao & \marcacao \\ \hline
                    08 & \marcacao & \marcacao & \marcacao & \marcacao \\ \hline
                    09 & \marcacao & \marcacao & \marcacao & \marcacao \\ \hline
                    10 & \marcacao & \marcacao & \marcacao & \marcacao \\ \hline
                    11 & \marcacao & \marcacao & \marcacao & \marcacao \\ \hline
                    12 & \marcacao & \marcacao & \marcacao & \marcacao \\ \hline
                    13 & \marcacao & \marcacao & \marcacao & \marcacao \\ \hline
                    14 & \marcacao & \marcacao & \marcacao & \marcacao \\ \hline
                    15 & \marcacao & \marcacao & \marcacao & \marcacao \\ \hline
                    16 & \marcacao & \marcacao & \marcacao & \marcacao \\ \hline
                    17 & \marcacao & \marcacao & \marcacao & \marcacao \\ \hline
                    18 & \marcacao & \marcacao & \marcacao & \marcacao \\ \hline
                    19 & \marcacao & \marcacao & \marcacao & \marcacao \\ \hline
                    20 & \marcacao & \marcacao & \marcacao & \marcacao \\ \hline
                    21 & \marcacao & \marcacao & \marcacao & \marcacao \\ \hline
                    22 & \marcacao & \marcacao & \marcacao & \marcacao \\ \hline
                    23 & \marcacao & \marcacao & \marcacao & \marcacao \\ \hline
                    24 & \marcacao & \marcacao & \marcacao & \marcacao \\ \hline
                    25 & \marcacao & \marcacao & \marcacao & \marcacao \\ \hline
                \end{tabular}
            };
            
            % --- Marcadores OMR ---
            \fill[black] ([shift={(-7mm,7mm)}]table.north west) rectangle ++(5mm,-5mm);
            \fill[black] ([shift={(7mm,7mm)}]table.north east) rectangle ++(-5mm,-5mm);
            \fill[black] ([shift={(-7mm,-7mm)}]table.south west) rectangle ++(5mm,5mm);
            \fill[black] ([shift={(7mm,-7mm)}]table.south east) rectangle ++(-5mm,5mm);
        \end{tikzpicture}
    \end{center}
    \renewcommand{\arraystretch}{1.0}
}

% Chamada do gabarito
\Gabarito

\newpage
% --- Bolinha centralizada ---
\newcommand{\marcacao}{%
    \vfill
    \tikz[baseline=(c.base)] \node (c) [circle, draw, minimum width=0.45cm, minimum height=0.45cm] {};
    \vfill
}

% --- Nova coluna centralizada ---
\newcolumntype{C}[1]{>{\centering\arraybackslash}m{#1}}

% --- Cabeçalho ---
\begin{minipage}[t]{0.7\textwidth}
    \textbf{Escola:} CEF04\_PLAN \\
    \textbf{Turma:} 2L - 9º ANO - Matutino \\
    \textbf{Aluno:} ISABELLA MOTA NASCIMENTO
\end{minipage}
\begin{minipage}[t]{0.29\textwidth}
    \raggedleft
    \qrcode[height=2.2cm]{ 573404-ISABELLA_MOTA_NASCIMENTO-2L }
\end{minipage}

\vspace{0.5cm}
\hrule
\vspace{0.5cm}

% --- Gabarito com 25 questões compactado ---
\newcommand{\Gabarito}{
    \renewcommand{\arraystretch}{1.2} % altura das linhas compacta
    \begin{center}
        \begin{tikzpicture}
            \node (table) [inner sep=0pt, anchor=north west] {
                \begin{tabular}{|C{0.6cm}|C{0.8cm}|C{0.8cm}|C{0.8cm}|C{0.8cm}|}
                    \hline
                    \adjustbox{max width=0.5cm}{\textbf{Nº}} & \textbf{A} & \textbf{B} & \textbf{C} & \textbf{D} \\
                    \hline
                    01 & \marcacao & \marcacao & \marcacao & \marcacao \\ \hline
                    02 & \marcacao & \marcacao & \marcacao & \marcacao \\ \hline
                    03 & \marcacao & \marcacao & \marcacao & \marcacao \\ \hline
                    04 & \marcacao & \marcacao & \marcacao & \marcacao \\ \hline
                    05 & \marcacao & \marcacao & \marcacao & \marcacao \\ \hline
                    06 & \marcacao & \marcacao & \marcacao & \marcacao \\ \hline
                    07 & \marcacao & \marcacao & \marcacao & \marcacao \\ \hline
                    08 & \marcacao & \marcacao & \marcacao & \marcacao \\ \hline
                    09 & \marcacao & \marcacao & \marcacao & \marcacao \\ \hline
                    10 & \marcacao & \marcacao & \marcacao & \marcacao \\ \hline
                    11 & \marcacao & \marcacao & \marcacao & \marcacao \\ \hline
                    12 & \marcacao & \marcacao & \marcacao & \marcacao \\ \hline
                    13 & \marcacao & \marcacao & \marcacao & \marcacao \\ \hline
                    14 & \marcacao & \marcacao & \marcacao & \marcacao \\ \hline
                    15 & \marcacao & \marcacao & \marcacao & \marcacao \\ \hline
                    16 & \marcacao & \marcacao & \marcacao & \marcacao \\ \hline
                    17 & \marcacao & \marcacao & \marcacao & \marcacao \\ \hline
                    18 & \marcacao & \marcacao & \marcacao & \marcacao \\ \hline
                    19 & \marcacao & \marcacao & \marcacao & \marcacao \\ \hline
                    20 & \marcacao & \marcacao & \marcacao & \marcacao \\ \hline
                    21 & \marcacao & \marcacao & \marcacao & \marcacao \\ \hline
                    22 & \marcacao & \marcacao & \marcacao & \marcacao \\ \hline
                    23 & \marcacao & \marcacao & \marcacao & \marcacao \\ \hline
                    24 & \marcacao & \marcacao & \marcacao & \marcacao \\ \hline
                    25 & \marcacao & \marcacao & \marcacao & \marcacao \\ \hline
                \end{tabular}
            };
            
            % --- Marcadores OMR ---
            \fill[black] ([shift={(-7mm,7mm)}]table.north west) rectangle ++(5mm,-5mm);
            \fill[black] ([shift={(7mm,7mm)}]table.north east) rectangle ++(-5mm,-5mm);
            \fill[black] ([shift={(-7mm,-7mm)}]table.south west) rectangle ++(5mm,5mm);
            \fill[black] ([shift={(7mm,-7mm)}]table.south east) rectangle ++(-5mm,5mm);
        \end{tikzpicture}
    \end{center}
    \renewcommand{\arraystretch}{1.0}
}

% Chamada do gabarito
\Gabarito

\newpage
% --- Bolinha centralizada ---
\newcommand{\marcacao}{%
    \vfill
    \tikz[baseline=(c.base)] \node (c) [circle, draw, minimum width=0.45cm, minimum height=0.45cm] {};
    \vfill
}

% --- Nova coluna centralizada ---
\newcolumntype{C}[1]{>{\centering\arraybackslash}m{#1}}

% --- Cabeçalho ---
\begin{minipage}[t]{0.7\textwidth}
    \textbf{Escola:} CEF04\_PLAN \\
    \textbf{Turma:} 2L - 9º ANO - Matutino \\
    \textbf{Aluno:} ÍTALO EMERSON DOS SANTOS SILVA
\end{minipage}
\begin{minipage}[t]{0.29\textwidth}
    \raggedleft
    \qrcode[height=2.2cm]{ 466372-ÍTALO_EMERSON_DOS_SANTOS_SILVA-2L }
\end{minipage}

\vspace{0.5cm}
\hrule
\vspace{0.5cm}

% --- Gabarito com 25 questões compactado ---
\newcommand{\Gabarito}{
    \renewcommand{\arraystretch}{1.2} % altura das linhas compacta
    \begin{center}
        \begin{tikzpicture}
            \node (table) [inner sep=0pt, anchor=north west] {
                \begin{tabular}{|C{0.6cm}|C{0.8cm}|C{0.8cm}|C{0.8cm}|C{0.8cm}|}
                    \hline
                    \adjustbox{max width=0.5cm}{\textbf{Nº}} & \textbf{A} & \textbf{B} & \textbf{C} & \textbf{D} \\
                    \hline
                    01 & \marcacao & \marcacao & \marcacao & \marcacao \\ \hline
                    02 & \marcacao & \marcacao & \marcacao & \marcacao \\ \hline
                    03 & \marcacao & \marcacao & \marcacao & \marcacao \\ \hline
                    04 & \marcacao & \marcacao & \marcacao & \marcacao \\ \hline
                    05 & \marcacao & \marcacao & \marcacao & \marcacao \\ \hline
                    06 & \marcacao & \marcacao & \marcacao & \marcacao \\ \hline
                    07 & \marcacao & \marcacao & \marcacao & \marcacao \\ \hline
                    08 & \marcacao & \marcacao & \marcacao & \marcacao \\ \hline
                    09 & \marcacao & \marcacao & \marcacao & \marcacao \\ \hline
                    10 & \marcacao & \marcacao & \marcacao & \marcacao \\ \hline
                    11 & \marcacao & \marcacao & \marcacao & \marcacao \\ \hline
                    12 & \marcacao & \marcacao & \marcacao & \marcacao \\ \hline
                    13 & \marcacao & \marcacao & \marcacao & \marcacao \\ \hline
                    14 & \marcacao & \marcacao & \marcacao & \marcacao \\ \hline
                    15 & \marcacao & \marcacao & \marcacao & \marcacao \\ \hline
                    16 & \marcacao & \marcacao & \marcacao & \marcacao \\ \hline
                    17 & \marcacao & \marcacao & \marcacao & \marcacao \\ \hline
                    18 & \marcacao & \marcacao & \marcacao & \marcacao \\ \hline
                    19 & \marcacao & \marcacao & \marcacao & \marcacao \\ \hline
                    20 & \marcacao & \marcacao & \marcacao & \marcacao \\ \hline
                    21 & \marcacao & \marcacao & \marcacao & \marcacao \\ \hline
                    22 & \marcacao & \marcacao & \marcacao & \marcacao \\ \hline
                    23 & \marcacao & \marcacao & \marcacao & \marcacao \\ \hline
                    24 & \marcacao & \marcacao & \marcacao & \marcacao \\ \hline
                    25 & \marcacao & \marcacao & \marcacao & \marcacao \\ \hline
                \end{tabular}
            };
            
            % --- Marcadores OMR ---
            \fill[black] ([shift={(-7mm,7mm)}]table.north west) rectangle ++(5mm,-5mm);
            \fill[black] ([shift={(7mm,7mm)}]table.north east) rectangle ++(-5mm,-5mm);
            \fill[black] ([shift={(-7mm,-7mm)}]table.south west) rectangle ++(5mm,5mm);
            \fill[black] ([shift={(7mm,-7mm)}]table.south east) rectangle ++(-5mm,5mm);
        \end{tikzpicture}
    \end{center}
    \renewcommand{\arraystretch}{1.0}
}

% Chamada do gabarito
\Gabarito

\newpage
% --- Bolinha centralizada ---
\newcommand{\marcacao}{%
    \vfill
    \tikz[baseline=(c.base)] \node (c) [circle, draw, minimum width=0.45cm, minimum height=0.45cm] {};
    \vfill
}

% --- Nova coluna centralizada ---
\newcolumntype{C}[1]{>{\centering\arraybackslash}m{#1}}

% --- Cabeçalho ---
\begin{minipage}[t]{0.7\textwidth}
    \textbf{Escola:} CEF04\_PLAN \\
    \textbf{Turma:} 2L - 9º ANO - Matutino \\
    \textbf{Aluno:} JOAO CARLOS RODRIGUES CARVALHO
\end{minipage}
\begin{minipage}[t]{0.29\textwidth}
    \raggedleft
    \qrcode[height=2.2cm]{ 560288-JOAO_CARLOS_RODRIGUES_CARVALHO-2L }
\end{minipage}

\vspace{0.5cm}
\hrule
\vspace{0.5cm}

% --- Gabarito com 25 questões compactado ---
\newcommand{\Gabarito}{
    \renewcommand{\arraystretch}{1.2} % altura das linhas compacta
    \begin{center}
        \begin{tikzpicture}
            \node (table) [inner sep=0pt, anchor=north west] {
                \begin{tabular}{|C{0.6cm}|C{0.8cm}|C{0.8cm}|C{0.8cm}|C{0.8cm}|}
                    \hline
                    \adjustbox{max width=0.5cm}{\textbf{Nº}} & \textbf{A} & \textbf{B} & \textbf{C} & \textbf{D} \\
                    \hline
                    01 & \marcacao & \marcacao & \marcacao & \marcacao \\ \hline
                    02 & \marcacao & \marcacao & \marcacao & \marcacao \\ \hline
                    03 & \marcacao & \marcacao & \marcacao & \marcacao \\ \hline
                    04 & \marcacao & \marcacao & \marcacao & \marcacao \\ \hline
                    05 & \marcacao & \marcacao & \marcacao & \marcacao \\ \hline
                    06 & \marcacao & \marcacao & \marcacao & \marcacao \\ \hline
                    07 & \marcacao & \marcacao & \marcacao & \marcacao \\ \hline
                    08 & \marcacao & \marcacao & \marcacao & \marcacao \\ \hline
                    09 & \marcacao & \marcacao & \marcacao & \marcacao \\ \hline
                    10 & \marcacao & \marcacao & \marcacao & \marcacao \\ \hline
                    11 & \marcacao & \marcacao & \marcacao & \marcacao \\ \hline
                    12 & \marcacao & \marcacao & \marcacao & \marcacao \\ \hline
                    13 & \marcacao & \marcacao & \marcacao & \marcacao \\ \hline
                    14 & \marcacao & \marcacao & \marcacao & \marcacao \\ \hline
                    15 & \marcacao & \marcacao & \marcacao & \marcacao \\ \hline
                    16 & \marcacao & \marcacao & \marcacao & \marcacao \\ \hline
                    17 & \marcacao & \marcacao & \marcacao & \marcacao \\ \hline
                    18 & \marcacao & \marcacao & \marcacao & \marcacao \\ \hline
                    19 & \marcacao & \marcacao & \marcacao & \marcacao \\ \hline
                    20 & \marcacao & \marcacao & \marcacao & \marcacao \\ \hline
                    21 & \marcacao & \marcacao & \marcacao & \marcacao \\ \hline
                    22 & \marcacao & \marcacao & \marcacao & \marcacao \\ \hline
                    23 & \marcacao & \marcacao & \marcacao & \marcacao \\ \hline
                    24 & \marcacao & \marcacao & \marcacao & \marcacao \\ \hline
                    25 & \marcacao & \marcacao & \marcacao & \marcacao \\ \hline
                \end{tabular}
            };
            
            % --- Marcadores OMR ---
            \fill[black] ([shift={(-7mm,7mm)}]table.north west) rectangle ++(5mm,-5mm);
            \fill[black] ([shift={(7mm,7mm)}]table.north east) rectangle ++(-5mm,-5mm);
            \fill[black] ([shift={(-7mm,-7mm)}]table.south west) rectangle ++(5mm,5mm);
            \fill[black] ([shift={(7mm,-7mm)}]table.south east) rectangle ++(-5mm,5mm);
        \end{tikzpicture}
    \end{center}
    \renewcommand{\arraystretch}{1.0}
}

% Chamada do gabarito
\Gabarito

\newpage
% --- Bolinha centralizada ---
\newcommand{\marcacao}{%
    \vfill
    \tikz[baseline=(c.base)] \node (c) [circle, draw, minimum width=0.45cm, minimum height=0.45cm] {};
    \vfill
}

% --- Nova coluna centralizada ---
\newcolumntype{C}[1]{>{\centering\arraybackslash}m{#1}}

% --- Cabeçalho ---
\begin{minipage}[t]{0.7\textwidth}
    \textbf{Escola:} CEF04\_PLAN \\
    \textbf{Turma:} 2L - 9º ANO - Matutino \\
    \textbf{Aluno:} JOAO GABRIEL FERREIRA DOS SANTOS
\end{minipage}
\begin{minipage}[t]{0.29\textwidth}
    \raggedleft
    \qrcode[height=2.2cm]{ 571195-JOAO_GABRIEL_FERREIRA_DOS_SANTOS-2L }
\end{minipage}

\vspace{0.5cm}
\hrule
\vspace{0.5cm}

% --- Gabarito com 25 questões compactado ---
\newcommand{\Gabarito}{
    \renewcommand{\arraystretch}{1.2} % altura das linhas compacta
    \begin{center}
        \begin{tikzpicture}
            \node (table) [inner sep=0pt, anchor=north west] {
                \begin{tabular}{|C{0.6cm}|C{0.8cm}|C{0.8cm}|C{0.8cm}|C{0.8cm}|}
                    \hline
                    \adjustbox{max width=0.5cm}{\textbf{Nº}} & \textbf{A} & \textbf{B} & \textbf{C} & \textbf{D} \\
                    \hline
                    01 & \marcacao & \marcacao & \marcacao & \marcacao \\ \hline
                    02 & \marcacao & \marcacao & \marcacao & \marcacao \\ \hline
                    03 & \marcacao & \marcacao & \marcacao & \marcacao \\ \hline
                    04 & \marcacao & \marcacao & \marcacao & \marcacao \\ \hline
                    05 & \marcacao & \marcacao & \marcacao & \marcacao \\ \hline
                    06 & \marcacao & \marcacao & \marcacao & \marcacao \\ \hline
                    07 & \marcacao & \marcacao & \marcacao & \marcacao \\ \hline
                    08 & \marcacao & \marcacao & \marcacao & \marcacao \\ \hline
                    09 & \marcacao & \marcacao & \marcacao & \marcacao \\ \hline
                    10 & \marcacao & \marcacao & \marcacao & \marcacao \\ \hline
                    11 & \marcacao & \marcacao & \marcacao & \marcacao \\ \hline
                    12 & \marcacao & \marcacao & \marcacao & \marcacao \\ \hline
                    13 & \marcacao & \marcacao & \marcacao & \marcacao \\ \hline
                    14 & \marcacao & \marcacao & \marcacao & \marcacao \\ \hline
                    15 & \marcacao & \marcacao & \marcacao & \marcacao \\ \hline
                    16 & \marcacao & \marcacao & \marcacao & \marcacao \\ \hline
                    17 & \marcacao & \marcacao & \marcacao & \marcacao \\ \hline
                    18 & \marcacao & \marcacao & \marcacao & \marcacao \\ \hline
                    19 & \marcacao & \marcacao & \marcacao & \marcacao \\ \hline
                    20 & \marcacao & \marcacao & \marcacao & \marcacao \\ \hline
                    21 & \marcacao & \marcacao & \marcacao & \marcacao \\ \hline
                    22 & \marcacao & \marcacao & \marcacao & \marcacao \\ \hline
                    23 & \marcacao & \marcacao & \marcacao & \marcacao \\ \hline
                    24 & \marcacao & \marcacao & \marcacao & \marcacao \\ \hline
                    25 & \marcacao & \marcacao & \marcacao & \marcacao \\ \hline
                \end{tabular}
            };
            
            % --- Marcadores OMR ---
            \fill[black] ([shift={(-7mm,7mm)}]table.north west) rectangle ++(5mm,-5mm);
            \fill[black] ([shift={(7mm,7mm)}]table.north east) rectangle ++(-5mm,-5mm);
            \fill[black] ([shift={(-7mm,-7mm)}]table.south west) rectangle ++(5mm,5mm);
            \fill[black] ([shift={(7mm,-7mm)}]table.south east) rectangle ++(-5mm,5mm);
        \end{tikzpicture}
    \end{center}
    \renewcommand{\arraystretch}{1.0}
}

% Chamada do gabarito
\Gabarito

\newpage
% --- Bolinha centralizada ---
\newcommand{\marcacao}{%
    \vfill
    \tikz[baseline=(c.base)] \node (c) [circle, draw, minimum width=0.45cm, minimum height=0.45cm] {};
    \vfill
}

% --- Nova coluna centralizada ---
\newcolumntype{C}[1]{>{\centering\arraybackslash}m{#1}}

% --- Cabeçalho ---
\begin{minipage}[t]{0.7\textwidth}
    \textbf{Escola:} CEF04\_PLAN \\
    \textbf{Turma:} 2L - 9º ANO - Matutino \\
    \textbf{Aluno:} PEDRO DAVI BARBOSA XAVIER
\end{minipage}
\begin{minipage}[t]{0.29\textwidth}
    \raggedleft
    \qrcode[height=2.2cm]{ 1071715-PEDRO_DAVI_BARBOSA_XAVIER-2L }
\end{minipage}

\vspace{0.5cm}
\hrule
\vspace{0.5cm}

% --- Gabarito com 25 questões compactado ---
\newcommand{\Gabarito}{
    \renewcommand{\arraystretch}{1.2} % altura das linhas compacta
    \begin{center}
        \begin{tikzpicture}
            \node (table) [inner sep=0pt, anchor=north west] {
                \begin{tabular}{|C{0.6cm}|C{0.8cm}|C{0.8cm}|C{0.8cm}|C{0.8cm}|}
                    \hline
                    \adjustbox{max width=0.5cm}{\textbf{Nº}} & \textbf{A} & \textbf{B} & \textbf{C} & \textbf{D} \\
                    \hline
                    01 & \marcacao & \marcacao & \marcacao & \marcacao \\ \hline
                    02 & \marcacao & \marcacao & \marcacao & \marcacao \\ \hline
                    03 & \marcacao & \marcacao & \marcacao & \marcacao \\ \hline
                    04 & \marcacao & \marcacao & \marcacao & \marcacao \\ \hline
                    05 & \marcacao & \marcacao & \marcacao & \marcacao \\ \hline
                    06 & \marcacao & \marcacao & \marcacao & \marcacao \\ \hline
                    07 & \marcacao & \marcacao & \marcacao & \marcacao \\ \hline
                    08 & \marcacao & \marcacao & \marcacao & \marcacao \\ \hline
                    09 & \marcacao & \marcacao & \marcacao & \marcacao \\ \hline
                    10 & \marcacao & \marcacao & \marcacao & \marcacao \\ \hline
                    11 & \marcacao & \marcacao & \marcacao & \marcacao \\ \hline
                    12 & \marcacao & \marcacao & \marcacao & \marcacao \\ \hline
                    13 & \marcacao & \marcacao & \marcacao & \marcacao \\ \hline
                    14 & \marcacao & \marcacao & \marcacao & \marcacao \\ \hline
                    15 & \marcacao & \marcacao & \marcacao & \marcacao \\ \hline
                    16 & \marcacao & \marcacao & \marcacao & \marcacao \\ \hline
                    17 & \marcacao & \marcacao & \marcacao & \marcacao \\ \hline
                    18 & \marcacao & \marcacao & \marcacao & \marcacao \\ \hline
                    19 & \marcacao & \marcacao & \marcacao & \marcacao \\ \hline
                    20 & \marcacao & \marcacao & \marcacao & \marcacao \\ \hline
                    21 & \marcacao & \marcacao & \marcacao & \marcacao \\ \hline
                    22 & \marcacao & \marcacao & \marcacao & \marcacao \\ \hline
                    23 & \marcacao & \marcacao & \marcacao & \marcacao \\ \hline
                    24 & \marcacao & \marcacao & \marcacao & \marcacao \\ \hline
                    25 & \marcacao & \marcacao & \marcacao & \marcacao \\ \hline
                \end{tabular}
            };
            
            % --- Marcadores OMR ---
            \fill[black] ([shift={(-7mm,7mm)}]table.north west) rectangle ++(5mm,-5mm);
            \fill[black] ([shift={(7mm,7mm)}]table.north east) rectangle ++(-5mm,-5mm);
            \fill[black] ([shift={(-7mm,-7mm)}]table.south west) rectangle ++(5mm,5mm);
            \fill[black] ([shift={(7mm,-7mm)}]table.south east) rectangle ++(-5mm,5mm);
        \end{tikzpicture}
    \end{center}
    \renewcommand{\arraystretch}{1.0}
}

% Chamada do gabarito
\Gabarito

\newpage
% --- Bolinha centralizada ---
\newcommand{\marcacao}{%
    \vfill
    \tikz[baseline=(c.base)] \node (c) [circle, draw, minimum width=0.45cm, minimum height=0.45cm] {};
    \vfill
}

% --- Nova coluna centralizada ---
\newcolumntype{C}[1]{>{\centering\arraybackslash}m{#1}}

% --- Cabeçalho ---
\begin{minipage}[t]{0.7\textwidth}
    \textbf{Escola:} CEF04\_PLAN \\
    \textbf{Turma:} 2L - 9º ANO - Matutino \\
    \textbf{Aluno:} SAMIRA NUNES ALMEIDA
\end{minipage}
\begin{minipage}[t]{0.29\textwidth}
    \raggedleft
    \qrcode[height=2.2cm]{ 1219509-SAMIRA_NUNES_ALMEIDA-2L }
\end{minipage}

\vspace{0.5cm}
\hrule
\vspace{0.5cm}

% --- Gabarito com 25 questões compactado ---
\newcommand{\Gabarito}{
    \renewcommand{\arraystretch}{1.2} % altura das linhas compacta
    \begin{center}
        \begin{tikzpicture}
            \node (table) [inner sep=0pt, anchor=north west] {
                \begin{tabular}{|C{0.6cm}|C{0.8cm}|C{0.8cm}|C{0.8cm}|C{0.8cm}|}
                    \hline
                    \adjustbox{max width=0.5cm}{\textbf{Nº}} & \textbf{A} & \textbf{B} & \textbf{C} & \textbf{D} \\
                    \hline
                    01 & \marcacao & \marcacao & \marcacao & \marcacao \\ \hline
                    02 & \marcacao & \marcacao & \marcacao & \marcacao \\ \hline
                    03 & \marcacao & \marcacao & \marcacao & \marcacao \\ \hline
                    04 & \marcacao & \marcacao & \marcacao & \marcacao \\ \hline
                    05 & \marcacao & \marcacao & \marcacao & \marcacao \\ \hline
                    06 & \marcacao & \marcacao & \marcacao & \marcacao \\ \hline
                    07 & \marcacao & \marcacao & \marcacao & \marcacao \\ \hline
                    08 & \marcacao & \marcacao & \marcacao & \marcacao \\ \hline
                    09 & \marcacao & \marcacao & \marcacao & \marcacao \\ \hline
                    10 & \marcacao & \marcacao & \marcacao & \marcacao \\ \hline
                    11 & \marcacao & \marcacao & \marcacao & \marcacao \\ \hline
                    12 & \marcacao & \marcacao & \marcacao & \marcacao \\ \hline
                    13 & \marcacao & \marcacao & \marcacao & \marcacao \\ \hline
                    14 & \marcacao & \marcacao & \marcacao & \marcacao \\ \hline
                    15 & \marcacao & \marcacao & \marcacao & \marcacao \\ \hline
                    16 & \marcacao & \marcacao & \marcacao & \marcacao \\ \hline
                    17 & \marcacao & \marcacao & \marcacao & \marcacao \\ \hline
                    18 & \marcacao & \marcacao & \marcacao & \marcacao \\ \hline
                    19 & \marcacao & \marcacao & \marcacao & \marcacao \\ \hline
                    20 & \marcacao & \marcacao & \marcacao & \marcacao \\ \hline
                    21 & \marcacao & \marcacao & \marcacao & \marcacao \\ \hline
                    22 & \marcacao & \marcacao & \marcacao & \marcacao \\ \hline
                    23 & \marcacao & \marcacao & \marcacao & \marcacao \\ \hline
                    24 & \marcacao & \marcacao & \marcacao & \marcacao \\ \hline
                    25 & \marcacao & \marcacao & \marcacao & \marcacao \\ \hline
                \end{tabular}
            };
            
            % --- Marcadores OMR ---
            \fill[black] ([shift={(-7mm,7mm)}]table.north west) rectangle ++(5mm,-5mm);
            \fill[black] ([shift={(7mm,7mm)}]table.north east) rectangle ++(-5mm,-5mm);
            \fill[black] ([shift={(-7mm,-7mm)}]table.south west) rectangle ++(5mm,5mm);
            \fill[black] ([shift={(7mm,-7mm)}]table.south east) rectangle ++(-5mm,5mm);
        \end{tikzpicture}
    \end{center}
    \renewcommand{\arraystretch}{1.0}
}

% Chamada do gabarito
\Gabarito

\newpage
% --- Bolinha centralizada ---
\newcommand{\marcacao}{%
    \vfill
    \tikz[baseline=(c.base)] \node (c) [circle, draw, minimum width=0.45cm, minimum height=0.45cm] {};
    \vfill
}

% --- Nova coluna centralizada ---
\newcolumntype{C}[1]{>{\centering\arraybackslash}m{#1}}

% --- Cabeçalho ---
\begin{minipage}[t]{0.7\textwidth}
    \textbf{Escola:} CEF04\_PLAN \\
    \textbf{Turma:} 2L - 9º ANO - Matutino \\
    \textbf{Aluno:} SAMUEL DAVI PAULINO MACIEL
\end{minipage}
\begin{minipage}[t]{0.29\textwidth}
    \raggedleft
    \qrcode[height=2.2cm]{ 592664-SAMUEL_DAVI_PAULINO_MACIEL-2L }
\end{minipage}

\vspace{0.5cm}
\hrule
\vspace{0.5cm}

% --- Gabarito com 25 questões compactado ---
\newcommand{\Gabarito}{
    \renewcommand{\arraystretch}{1.2} % altura das linhas compacta
    \begin{center}
        \begin{tikzpicture}
            \node (table) [inner sep=0pt, anchor=north west] {
                \begin{tabular}{|C{0.6cm}|C{0.8cm}|C{0.8cm}|C{0.8cm}|C{0.8cm}|}
                    \hline
                    \adjustbox{max width=0.5cm}{\textbf{Nº}} & \textbf{A} & \textbf{B} & \textbf{C} & \textbf{D} \\
                    \hline
                    01 & \marcacao & \marcacao & \marcacao & \marcacao \\ \hline
                    02 & \marcacao & \marcacao & \marcacao & \marcacao \\ \hline
                    03 & \marcacao & \marcacao & \marcacao & \marcacao \\ \hline
                    04 & \marcacao & \marcacao & \marcacao & \marcacao \\ \hline
                    05 & \marcacao & \marcacao & \marcacao & \marcacao \\ \hline
                    06 & \marcacao & \marcacao & \marcacao & \marcacao \\ \hline
                    07 & \marcacao & \marcacao & \marcacao & \marcacao \\ \hline
                    08 & \marcacao & \marcacao & \marcacao & \marcacao \\ \hline
                    09 & \marcacao & \marcacao & \marcacao & \marcacao \\ \hline
                    10 & \marcacao & \marcacao & \marcacao & \marcacao \\ \hline
                    11 & \marcacao & \marcacao & \marcacao & \marcacao \\ \hline
                    12 & \marcacao & \marcacao & \marcacao & \marcacao \\ \hline
                    13 & \marcacao & \marcacao & \marcacao & \marcacao \\ \hline
                    14 & \marcacao & \marcacao & \marcacao & \marcacao \\ \hline
                    15 & \marcacao & \marcacao & \marcacao & \marcacao \\ \hline
                    16 & \marcacao & \marcacao & \marcacao & \marcacao \\ \hline
                    17 & \marcacao & \marcacao & \marcacao & \marcacao \\ \hline
                    18 & \marcacao & \marcacao & \marcacao & \marcacao \\ \hline
                    19 & \marcacao & \marcacao & \marcacao & \marcacao \\ \hline
                    20 & \marcacao & \marcacao & \marcacao & \marcacao \\ \hline
                    21 & \marcacao & \marcacao & \marcacao & \marcacao \\ \hline
                    22 & \marcacao & \marcacao & \marcacao & \marcacao \\ \hline
                    23 & \marcacao & \marcacao & \marcacao & \marcacao \\ \hline
                    24 & \marcacao & \marcacao & \marcacao & \marcacao \\ \hline
                    25 & \marcacao & \marcacao & \marcacao & \marcacao \\ \hline
                \end{tabular}
            };
            
            % --- Marcadores OMR ---
            \fill[black] ([shift={(-7mm,7mm)}]table.north west) rectangle ++(5mm,-5mm);
            \fill[black] ([shift={(7mm,7mm)}]table.north east) rectangle ++(-5mm,-5mm);
            \fill[black] ([shift={(-7mm,-7mm)}]table.south west) rectangle ++(5mm,5mm);
            \fill[black] ([shift={(7mm,-7mm)}]table.south east) rectangle ++(-5mm,5mm);
        \end{tikzpicture}
    \end{center}
    \renewcommand{\arraystretch}{1.0}
}

% Chamada do gabarito
\Gabarito

\newpage
% --- Bolinha centralizada ---
\newcommand{\marcacao}{%
    \vfill
    \tikz[baseline=(c.base)] \node (c) [circle, draw, minimum width=0.45cm, minimum height=0.45cm] {};
    \vfill
}

% --- Nova coluna centralizada ---
\newcolumntype{C}[1]{>{\centering\arraybackslash}m{#1}}

% --- Cabeçalho ---
\begin{minipage}[t]{0.7\textwidth}
    \textbf{Escola:} CEF04\_PLAN \\
    \textbf{Turma:} 2L - 9º ANO - Matutino \\
    \textbf{Aluno:} SAMUEL DE OLIVEIRA SOARES
\end{minipage}
\begin{minipage}[t]{0.29\textwidth}
    \raggedleft
    \qrcode[height=2.2cm]{ 532862-SAMUEL_DE_OLIVEIRA_SOARES-2L }
\end{minipage}

\vspace{0.5cm}
\hrule
\vspace{0.5cm}

% --- Gabarito com 25 questões compactado ---
\newcommand{\Gabarito}{
    \renewcommand{\arraystretch}{1.2} % altura das linhas compacta
    \begin{center}
        \begin{tikzpicture}
            \node (table) [inner sep=0pt, anchor=north west] {
                \begin{tabular}{|C{0.6cm}|C{0.8cm}|C{0.8cm}|C{0.8cm}|C{0.8cm}|}
                    \hline
                    \adjustbox{max width=0.5cm}{\textbf{Nº}} & \textbf{A} & \textbf{B} & \textbf{C} & \textbf{D} \\
                    \hline
                    01 & \marcacao & \marcacao & \marcacao & \marcacao \\ \hline
                    02 & \marcacao & \marcacao & \marcacao & \marcacao \\ \hline
                    03 & \marcacao & \marcacao & \marcacao & \marcacao \\ \hline
                    04 & \marcacao & \marcacao & \marcacao & \marcacao \\ \hline
                    05 & \marcacao & \marcacao & \marcacao & \marcacao \\ \hline
                    06 & \marcacao & \marcacao & \marcacao & \marcacao \\ \hline
                    07 & \marcacao & \marcacao & \marcacao & \marcacao \\ \hline
                    08 & \marcacao & \marcacao & \marcacao & \marcacao \\ \hline
                    09 & \marcacao & \marcacao & \marcacao & \marcacao \\ \hline
                    10 & \marcacao & \marcacao & \marcacao & \marcacao \\ \hline
                    11 & \marcacao & \marcacao & \marcacao & \marcacao \\ \hline
                    12 & \marcacao & \marcacao & \marcacao & \marcacao \\ \hline
                    13 & \marcacao & \marcacao & \marcacao & \marcacao \\ \hline
                    14 & \marcacao & \marcacao & \marcacao & \marcacao \\ \hline
                    15 & \marcacao & \marcacao & \marcacao & \marcacao \\ \hline
                    16 & \marcacao & \marcacao & \marcacao & \marcacao \\ \hline
                    17 & \marcacao & \marcacao & \marcacao & \marcacao \\ \hline
                    18 & \marcacao & \marcacao & \marcacao & \marcacao \\ \hline
                    19 & \marcacao & \marcacao & \marcacao & \marcacao \\ \hline
                    20 & \marcacao & \marcacao & \marcacao & \marcacao \\ \hline
                    21 & \marcacao & \marcacao & \marcacao & \marcacao \\ \hline
                    22 & \marcacao & \marcacao & \marcacao & \marcacao \\ \hline
                    23 & \marcacao & \marcacao & \marcacao & \marcacao \\ \hline
                    24 & \marcacao & \marcacao & \marcacao & \marcacao \\ \hline
                    25 & \marcacao & \marcacao & \marcacao & \marcacao \\ \hline
                \end{tabular}
            };
            
            % --- Marcadores OMR ---
            \fill[black] ([shift={(-7mm,7mm)}]table.north west) rectangle ++(5mm,-5mm);
            \fill[black] ([shift={(7mm,7mm)}]table.north east) rectangle ++(-5mm,-5mm);
            \fill[black] ([shift={(-7mm,-7mm)}]table.south west) rectangle ++(5mm,5mm);
            \fill[black] ([shift={(7mm,-7mm)}]table.south east) rectangle ++(-5mm,5mm);
        \end{tikzpicture}
    \end{center}
    \renewcommand{\arraystretch}{1.0}
}

% Chamada do gabarito
\Gabarito

\newpage
% --- Bolinha centralizada ---
\newcommand{\marcacao}{%
    \vfill
    \tikz[baseline=(c.base)] \node (c) [circle, draw, minimum width=0.45cm, minimum height=0.45cm] {};
    \vfill
}

% --- Nova coluna centralizada ---
\newcolumntype{C}[1]{>{\centering\arraybackslash}m{#1}}

% --- Cabeçalho ---
\begin{minipage}[t]{0.7\textwidth}
    \textbf{Escola:} CEF04\_PLAN \\
    \textbf{Turma:} 2L - 9º ANO - Matutino \\
    \textbf{Aluno:} THAYLA YASMIN ALVES DOS SANTOS
\end{minipage}
\begin{minipage}[t]{0.29\textwidth}
    \raggedleft
    \qrcode[height=2.2cm]{ 468742-THAYLA_YASMIN_ALVES_DOS_SANTOS-2L }
\end{minipage}

\vspace{0.5cm}
\hrule
\vspace{0.5cm}

% --- Gabarito com 25 questões compactado ---
\newcommand{\Gabarito}{
    \renewcommand{\arraystretch}{1.2} % altura das linhas compacta
    \begin{center}
        \begin{tikzpicture}
            \node (table) [inner sep=0pt, anchor=north west] {
                \begin{tabular}{|C{0.6cm}|C{0.8cm}|C{0.8cm}|C{0.8cm}|C{0.8cm}|}
                    \hline
                    \adjustbox{max width=0.5cm}{\textbf{Nº}} & \textbf{A} & \textbf{B} & \textbf{C} & \textbf{D} \\
                    \hline
                    01 & \marcacao & \marcacao & \marcacao & \marcacao \\ \hline
                    02 & \marcacao & \marcacao & \marcacao & \marcacao \\ \hline
                    03 & \marcacao & \marcacao & \marcacao & \marcacao \\ \hline
                    04 & \marcacao & \marcacao & \marcacao & \marcacao \\ \hline
                    05 & \marcacao & \marcacao & \marcacao & \marcacao \\ \hline
                    06 & \marcacao & \marcacao & \marcacao & \marcacao \\ \hline
                    07 & \marcacao & \marcacao & \marcacao & \marcacao \\ \hline
                    08 & \marcacao & \marcacao & \marcacao & \marcacao \\ \hline
                    09 & \marcacao & \marcacao & \marcacao & \marcacao \\ \hline
                    10 & \marcacao & \marcacao & \marcacao & \marcacao \\ \hline
                    11 & \marcacao & \marcacao & \marcacao & \marcacao \\ \hline
                    12 & \marcacao & \marcacao & \marcacao & \marcacao \\ \hline
                    13 & \marcacao & \marcacao & \marcacao & \marcacao \\ \hline
                    14 & \marcacao & \marcacao & \marcacao & \marcacao \\ \hline
                    15 & \marcacao & \marcacao & \marcacao & \marcacao \\ \hline
                    16 & \marcacao & \marcacao & \marcacao & \marcacao \\ \hline
                    17 & \marcacao & \marcacao & \marcacao & \marcacao \\ \hline
                    18 & \marcacao & \marcacao & \marcacao & \marcacao \\ \hline
                    19 & \marcacao & \marcacao & \marcacao & \marcacao \\ \hline
                    20 & \marcacao & \marcacao & \marcacao & \marcacao \\ \hline
                    21 & \marcacao & \marcacao & \marcacao & \marcacao \\ \hline
                    22 & \marcacao & \marcacao & \marcacao & \marcacao \\ \hline
                    23 & \marcacao & \marcacao & \marcacao & \marcacao \\ \hline
                    24 & \marcacao & \marcacao & \marcacao & \marcacao \\ \hline
                    25 & \marcacao & \marcacao & \marcacao & \marcacao \\ \hline
                \end{tabular}
            };
            
            % --- Marcadores OMR ---
            \fill[black] ([shift={(-7mm,7mm)}]table.north west) rectangle ++(5mm,-5mm);
            \fill[black] ([shift={(7mm,7mm)}]table.north east) rectangle ++(-5mm,-5mm);
            \fill[black] ([shift={(-7mm,-7mm)}]table.south west) rectangle ++(5mm,5mm);
            \fill[black] ([shift={(7mm,-7mm)}]table.south east) rectangle ++(-5mm,5mm);
        \end{tikzpicture}
    \end{center}
    \renewcommand{\arraystretch}{1.0}
}

% Chamada do gabarito
\Gabarito

\newpage
% --- Bolinha centralizada ---
\newcommand{\marcacao}{%
    \vfill
    \tikz[baseline=(c.base)] \node (c) [circle, draw, minimum width=0.45cm, minimum height=0.45cm] {};
    \vfill
}

% --- Nova coluna centralizada ---
\newcolumntype{C}[1]{>{\centering\arraybackslash}m{#1}}

% --- Cabeçalho ---
\begin{minipage}[t]{0.7\textwidth}
    \textbf{Escola:} CEF04\_PLAN \\
    \textbf{Turma:} 2L - 9º ANO - Matutino \\
    \textbf{Aluno:} VICTOR JUAN ANDRADE DOS SANTOS
\end{minipage}
\begin{minipage}[t]{0.29\textwidth}
    \raggedleft
    \qrcode[height=2.2cm]{ 540853-VICTOR_JUAN_ANDRADE_DOS_SANTOS-2L }
\end{minipage}

\vspace{0.5cm}
\hrule
\vspace{0.5cm}

% --- Gabarito com 25 questões compactado ---
\newcommand{\Gabarito}{
    \renewcommand{\arraystretch}{1.2} % altura das linhas compacta
    \begin{center}
        \begin{tikzpicture}
            \node (table) [inner sep=0pt, anchor=north west] {
                \begin{tabular}{|C{0.6cm}|C{0.8cm}|C{0.8cm}|C{0.8cm}|C{0.8cm}|}
                    \hline
                    \adjustbox{max width=0.5cm}{\textbf{Nº}} & \textbf{A} & \textbf{B} & \textbf{C} & \textbf{D} \\
                    \hline
                    01 & \marcacao & \marcacao & \marcacao & \marcacao \\ \hline
                    02 & \marcacao & \marcacao & \marcacao & \marcacao \\ \hline
                    03 & \marcacao & \marcacao & \marcacao & \marcacao \\ \hline
                    04 & \marcacao & \marcacao & \marcacao & \marcacao \\ \hline
                    05 & \marcacao & \marcacao & \marcacao & \marcacao \\ \hline
                    06 & \marcacao & \marcacao & \marcacao & \marcacao \\ \hline
                    07 & \marcacao & \marcacao & \marcacao & \marcacao \\ \hline
                    08 & \marcacao & \marcacao & \marcacao & \marcacao \\ \hline
                    09 & \marcacao & \marcacao & \marcacao & \marcacao \\ \hline
                    10 & \marcacao & \marcacao & \marcacao & \marcacao \\ \hline
                    11 & \marcacao & \marcacao & \marcacao & \marcacao \\ \hline
                    12 & \marcacao & \marcacao & \marcacao & \marcacao \\ \hline
                    13 & \marcacao & \marcacao & \marcacao & \marcacao \\ \hline
                    14 & \marcacao & \marcacao & \marcacao & \marcacao \\ \hline
                    15 & \marcacao & \marcacao & \marcacao & \marcacao \\ \hline
                    16 & \marcacao & \marcacao & \marcacao & \marcacao \\ \hline
                    17 & \marcacao & \marcacao & \marcacao & \marcacao \\ \hline
                    18 & \marcacao & \marcacao & \marcacao & \marcacao \\ \hline
                    19 & \marcacao & \marcacao & \marcacao & \marcacao \\ \hline
                    20 & \marcacao & \marcacao & \marcacao & \marcacao \\ \hline
                    21 & \marcacao & \marcacao & \marcacao & \marcacao \\ \hline
                    22 & \marcacao & \marcacao & \marcacao & \marcacao \\ \hline
                    23 & \marcacao & \marcacao & \marcacao & \marcacao \\ \hline
                    24 & \marcacao & \marcacao & \marcacao & \marcacao \\ \hline
                    25 & \marcacao & \marcacao & \marcacao & \marcacao \\ \hline
                \end{tabular}
            };
            
            % --- Marcadores OMR ---
            \fill[black] ([shift={(-7mm,7mm)}]table.north west) rectangle ++(5mm,-5mm);
            \fill[black] ([shift={(7mm,7mm)}]table.north east) rectangle ++(-5mm,-5mm);
            \fill[black] ([shift={(-7mm,-7mm)}]table.south west) rectangle ++(5mm,5mm);
            \fill[black] ([shift={(7mm,-7mm)}]table.south east) rectangle ++(-5mm,5mm);
        \end{tikzpicture}
    \end{center}
    \renewcommand{\arraystretch}{1.0}
}

% Chamada do gabarito
\Gabarito

\newpage
% --- Bolinha centralizada ---
\newcommand{\marcacao}{%
    \vfill
    \tikz[baseline=(c.base)] \node (c) [circle, draw, minimum width=0.45cm, minimum height=0.45cm] {};
    \vfill
}

% --- Nova coluna centralizada ---
\newcolumntype{C}[1]{>{\centering\arraybackslash}m{#1}}

% --- Cabeçalho ---
\begin{minipage}[t]{0.7\textwidth}
    \textbf{Escola:} CEF04\_PLAN \\
    \textbf{Turma:} 2L - 9º ANO - Matutino \\
    \textbf{Aluno:} VICTORIA MAGALHÃES CHAVES CAVALCANTI AGUIAR
\end{minipage}
\begin{minipage}[t]{0.29\textwidth}
    \raggedleft
    \qrcode[height=2.2cm]{ 1235947-VICTORIA_MAGALHÃES_CHAVES_CAVALCANTI_AGUIAR-2L }
\end{minipage}

\vspace{0.5cm}
\hrule
\vspace{0.5cm}

% --- Gabarito com 25 questões compactado ---
\newcommand{\Gabarito}{
    \renewcommand{\arraystretch}{1.2} % altura das linhas compacta
    \begin{center}
        \begin{tikzpicture}
            \node (table) [inner sep=0pt, anchor=north west] {
                \begin{tabular}{|C{0.6cm}|C{0.8cm}|C{0.8cm}|C{0.8cm}|C{0.8cm}|}
                    \hline
                    \adjustbox{max width=0.5cm}{\textbf{Nº}} & \textbf{A} & \textbf{B} & \textbf{C} & \textbf{D} \\
                    \hline
                    01 & \marcacao & \marcacao & \marcacao & \marcacao \\ \hline
                    02 & \marcacao & \marcacao & \marcacao & \marcacao \\ \hline
                    03 & \marcacao & \marcacao & \marcacao & \marcacao \\ \hline
                    04 & \marcacao & \marcacao & \marcacao & \marcacao \\ \hline
                    05 & \marcacao & \marcacao & \marcacao & \marcacao \\ \hline
                    06 & \marcacao & \marcacao & \marcacao & \marcacao \\ \hline
                    07 & \marcacao & \marcacao & \marcacao & \marcacao \\ \hline
                    08 & \marcacao & \marcacao & \marcacao & \marcacao \\ \hline
                    09 & \marcacao & \marcacao & \marcacao & \marcacao \\ \hline
                    10 & \marcacao & \marcacao & \marcacao & \marcacao \\ \hline
                    11 & \marcacao & \marcacao & \marcacao & \marcacao \\ \hline
                    12 & \marcacao & \marcacao & \marcacao & \marcacao \\ \hline
                    13 & \marcacao & \marcacao & \marcacao & \marcacao \\ \hline
                    14 & \marcacao & \marcacao & \marcacao & \marcacao \\ \hline
                    15 & \marcacao & \marcacao & \marcacao & \marcacao \\ \hline
                    16 & \marcacao & \marcacao & \marcacao & \marcacao \\ \hline
                    17 & \marcacao & \marcacao & \marcacao & \marcacao \\ \hline
                    18 & \marcacao & \marcacao & \marcacao & \marcacao \\ \hline
                    19 & \marcacao & \marcacao & \marcacao & \marcacao \\ \hline
                    20 & \marcacao & \marcacao & \marcacao & \marcacao \\ \hline
                    21 & \marcacao & \marcacao & \marcacao & \marcacao \\ \hline
                    22 & \marcacao & \marcacao & \marcacao & \marcacao \\ \hline
                    23 & \marcacao & \marcacao & \marcacao & \marcacao \\ \hline
                    24 & \marcacao & \marcacao & \marcacao & \marcacao \\ \hline
                    25 & \marcacao & \marcacao & \marcacao & \marcacao \\ \hline
                \end{tabular}
            };
            
            % --- Marcadores OMR ---
            \fill[black] ([shift={(-7mm,7mm)}]table.north west) rectangle ++(5mm,-5mm);
            \fill[black] ([shift={(7mm,7mm)}]table.north east) rectangle ++(-5mm,-5mm);
            \fill[black] ([shift={(-7mm,-7mm)}]table.south west) rectangle ++(5mm,5mm);
            \fill[black] ([shift={(7mm,-7mm)}]table.south east) rectangle ++(-5mm,5mm);
        \end{tikzpicture}
    \end{center}
    \renewcommand{\arraystretch}{1.0}
}

% Chamada do gabarito
\Gabarito

\newpage
% --- Bolinha centralizada ---
\newcommand{\marcacao}{%
    \vfill
    \tikz[baseline=(c.base)] \node (c) [circle, draw, minimum width=0.45cm, minimum height=0.45cm] {};
    \vfill
}

% --- Nova coluna centralizada ---
\newcolumntype{C}[1]{>{\centering\arraybackslash}m{#1}}

% --- Cabeçalho ---
\begin{minipage}[t]{0.7\textwidth}
    \textbf{Escola:} CEF04\_PLAN \\
    \textbf{Turma:} 2L - 9º ANO - Matutino \\
    \textbf{Aluno:} VITORIA RODRIGUES DE SOUSA
\end{minipage}
\begin{minipage}[t]{0.29\textwidth}
    \raggedleft
    \qrcode[height=2.2cm]{ 192101-VITORIA_RODRIGUES_DE_SOUSA-2L }
\end{minipage}

\vspace{0.5cm}
\hrule
\vspace{0.5cm}

% --- Gabarito com 25 questões compactado ---
\newcommand{\Gabarito}{
    \renewcommand{\arraystretch}{1.2} % altura das linhas compacta
    \begin{center}
        \begin{tikzpicture}
            \node (table) [inner sep=0pt, anchor=north west] {
                \begin{tabular}{|C{0.6cm}|C{0.8cm}|C{0.8cm}|C{0.8cm}|C{0.8cm}|}
                    \hline
                    \adjustbox{max width=0.5cm}{\textbf{Nº}} & \textbf{A} & \textbf{B} & \textbf{C} & \textbf{D} \\
                    \hline
                    01 & \marcacao & \marcacao & \marcacao & \marcacao \\ \hline
                    02 & \marcacao & \marcacao & \marcacao & \marcacao \\ \hline
                    03 & \marcacao & \marcacao & \marcacao & \marcacao \\ \hline
                    04 & \marcacao & \marcacao & \marcacao & \marcacao \\ \hline
                    05 & \marcacao & \marcacao & \marcacao & \marcacao \\ \hline
                    06 & \marcacao & \marcacao & \marcacao & \marcacao \\ \hline
                    07 & \marcacao & \marcacao & \marcacao & \marcacao \\ \hline
                    08 & \marcacao & \marcacao & \marcacao & \marcacao \\ \hline
                    09 & \marcacao & \marcacao & \marcacao & \marcacao \\ \hline
                    10 & \marcacao & \marcacao & \marcacao & \marcacao \\ \hline
                    11 & \marcacao & \marcacao & \marcacao & \marcacao \\ \hline
                    12 & \marcacao & \marcacao & \marcacao & \marcacao \\ \hline
                    13 & \marcacao & \marcacao & \marcacao & \marcacao \\ \hline
                    14 & \marcacao & \marcacao & \marcacao & \marcacao \\ \hline
                    15 & \marcacao & \marcacao & \marcacao & \marcacao \\ \hline
                    16 & \marcacao & \marcacao & \marcacao & \marcacao \\ \hline
                    17 & \marcacao & \marcacao & \marcacao & \marcacao \\ \hline
                    18 & \marcacao & \marcacao & \marcacao & \marcacao \\ \hline
                    19 & \marcacao & \marcacao & \marcacao & \marcacao \\ \hline
                    20 & \marcacao & \marcacao & \marcacao & \marcacao \\ \hline
                    21 & \marcacao & \marcacao & \marcacao & \marcacao \\ \hline
                    22 & \marcacao & \marcacao & \marcacao & \marcacao \\ \hline
                    23 & \marcacao & \marcacao & \marcacao & \marcacao \\ \hline
                    24 & \marcacao & \marcacao & \marcacao & \marcacao \\ \hline
                    25 & \marcacao & \marcacao & \marcacao & \marcacao \\ \hline
                \end{tabular}
            };
            
            % --- Marcadores OMR ---
            \fill[black] ([shift={(-7mm,7mm)}]table.north west) rectangle ++(5mm,-5mm);
            \fill[black] ([shift={(7mm,7mm)}]table.north east) rectangle ++(-5mm,-5mm);
            \fill[black] ([shift={(-7mm,-7mm)}]table.south west) rectangle ++(5mm,5mm);
            \fill[black] ([shift={(7mm,-7mm)}]table.south east) rectangle ++(-5mm,5mm);
        \end{tikzpicture}
    \end{center}
    \renewcommand{\arraystretch}{1.0}
}

% Chamada do gabarito
\Gabarito

\newpage
% --- Bolinha centralizada ---
\newcommand{\marcacao}{%
    \vfill
    \tikz[baseline=(c.base)] \node (c) [circle, draw, minimum width=0.45cm, minimum height=0.45cm] {};
    \vfill
}

% --- Nova coluna centralizada ---
\newcolumntype{C}[1]{>{\centering\arraybackslash}m{#1}}

% --- Cabeçalho ---
\begin{minipage}[t]{0.7\textwidth}
    \textbf{Escola:} CEF04\_PLAN \\
    \textbf{Turma:} 2L - 9º ANO - Matutino \\
    \textbf{Aluno:} YASMIN VITORIA CONDE DOS SANTOS
\end{minipage}
\begin{minipage}[t]{0.29\textwidth}
    \raggedleft
    \qrcode[height=2.2cm]{ 324619-YASMIN_VITORIA_CONDE_DOS_SANTOS-2L }
\end{minipage}

\vspace{0.5cm}
\hrule
\vspace{0.5cm}

% --- Gabarito com 25 questões compactado ---
\newcommand{\Gabarito}{
    \renewcommand{\arraystretch}{1.2} % altura das linhas compacta
    \begin{center}
        \begin{tikzpicture}
            \node (table) [inner sep=0pt, anchor=north west] {
                \begin{tabular}{|C{0.6cm}|C{0.8cm}|C{0.8cm}|C{0.8cm}|C{0.8cm}|}
                    \hline
                    \adjustbox{max width=0.5cm}{\textbf{Nº}} & \textbf{A} & \textbf{B} & \textbf{C} & \textbf{D} \\
                    \hline
                    01 & \marcacao & \marcacao & \marcacao & \marcacao \\ \hline
                    02 & \marcacao & \marcacao & \marcacao & \marcacao \\ \hline
                    03 & \marcacao & \marcacao & \marcacao & \marcacao \\ \hline
                    04 & \marcacao & \marcacao & \marcacao & \marcacao \\ \hline
                    05 & \marcacao & \marcacao & \marcacao & \marcacao \\ \hline
                    06 & \marcacao & \marcacao & \marcacao & \marcacao \\ \hline
                    07 & \marcacao & \marcacao & \marcacao & \marcacao \\ \hline
                    08 & \marcacao & \marcacao & \marcacao & \marcacao \\ \hline
                    09 & \marcacao & \marcacao & \marcacao & \marcacao \\ \hline
                    10 & \marcacao & \marcacao & \marcacao & \marcacao \\ \hline
                    11 & \marcacao & \marcacao & \marcacao & \marcacao \\ \hline
                    12 & \marcacao & \marcacao & \marcacao & \marcacao \\ \hline
                    13 & \marcacao & \marcacao & \marcacao & \marcacao \\ \hline
                    14 & \marcacao & \marcacao & \marcacao & \marcacao \\ \hline
                    15 & \marcacao & \marcacao & \marcacao & \marcacao \\ \hline
                    16 & \marcacao & \marcacao & \marcacao & \marcacao \\ \hline
                    17 & \marcacao & \marcacao & \marcacao & \marcacao \\ \hline
                    18 & \marcacao & \marcacao & \marcacao & \marcacao \\ \hline
                    19 & \marcacao & \marcacao & \marcacao & \marcacao \\ \hline
                    20 & \marcacao & \marcacao & \marcacao & \marcacao \\ \hline
                    21 & \marcacao & \marcacao & \marcacao & \marcacao \\ \hline
                    22 & \marcacao & \marcacao & \marcacao & \marcacao \\ \hline
                    23 & \marcacao & \marcacao & \marcacao & \marcacao \\ \hline
                    24 & \marcacao & \marcacao & \marcacao & \marcacao \\ \hline
                    25 & \marcacao & \marcacao & \marcacao & \marcacao \\ \hline
                \end{tabular}
            };
            
            % --- Marcadores OMR ---
            \fill[black] ([shift={(-7mm,7mm)}]table.north west) rectangle ++(5mm,-5mm);
            \fill[black] ([shift={(7mm,7mm)}]table.north east) rectangle ++(-5mm,-5mm);
            \fill[black] ([shift={(-7mm,-7mm)}]table.south west) rectangle ++(5mm,5mm);
            \fill[black] ([shift={(7mm,-7mm)}]table.south east) rectangle ++(-5mm,5mm);
        \end{tikzpicture}
    \end{center}
    \renewcommand{\arraystretch}{1.0}
}

% Chamada do gabarito
\Gabarito

\newpage
\end{document}